\documentclass[UTF8,11pt]{ctexart}
\usepackage{amsmath,amssymb,geometry}
\geometry{margin=22mm}
\begin{document}
\section*{P59: 长期平均收益最大化的动态随机匹配主问题}
\paragraph{来源与范围}依据原论文PDF第1, 6, 7, 8, 9, 10, 11, 20, 21, 22, 23页，第2节，式(1)及第9--10页政策/状态定义；第5节第20--21页$R_t$与长期平均收益；第21--22页(12)--(15)用于辨别基准模型。

\paragraph{保留的主问题}
采用原论文第5节的长期平均收益最大化动态匹配问题，沿用第2节的到达、
耐心与信息结构。下面以第10页所述的类型计数状态完整展开连续时间决策过程。


\paragraph{输入与随机机制}
类型集合 $\mathcal T=\{1,\ldots,d\}$ 有限。类型 $i$ 的参与者按速率
$\lambda_i>0$ 的泊松过程到达，各类型到达过程相互独立；每名类型 $i$
参与者的耐心为独立指数随机变量，参数 $0<\mu_i<\infty$，与到达过程独立。
参与者未匹配且耐心耗尽时立即离开。到达时间、类型与已经发生的离开可被观察，
未来到达及仍在系统者的耐心截止时间不可预知。

无向类型图上的匹配收益 $r_{ij}=r_{ji}\in\mathbb R$ 为给定有限数值。
同类型收益 $r_{ii}$ 对应两名不同的类型 $i$ 参与者。
可用零收益补全类型图（原文第20页）；没有放弃惩罚或等待成本。
双方参与者采用给定的有限耐心分布。

\paragraph{状态与完整动作域}
在一次外部事件发生之后、匹配决策之前，状态为
$\nu=(\nu_i)_{i\in\mathcal T}\in\mathbb Z_+^d$，表示当前仍可匹配的各类型人数。
动作 $a=(a_{ij})_{1\le i\le j\le d}$ 表示在该时点执行的各类无序匹配对数。
可行域为
\begin{equation}
\mathcal A(\nu)=\left\{a\in\mathbb Z_+^{d(d+1)/2}:\quad
2a_{ii}+\sum_{j<i}a_{ji}+\sum_{j>i}a_{ij}\le\nu_i,
\quad i=1,\ldots,d\right\}.\tag{A}
\end{equation}
允许空动作 $a=0$，即继续等待。每名参与者至多匹配一次，匹配双方立即从
可匹配池删除。匹配后剩余状态和本次收益分别为
\begin{align}
v_i(\nu,a)&=\nu_i-2a_{ii}-\sum_{j<i}a_{ji}-\sum_{j>i}a_{ij},\quad i\in\mathcal T,\tag{S}\\
g(a)&=\sum_{1\le i\le j\le d}r_{ij}a_{ij}.\tag{R}
\end{align}
式(A)是原文“在当前剩余顶点诱导子图上选匹配”的类型计数形式；
各配对数均为整数且满足人数容量时，可在该类型完全图中选择相互不重叠的实际参与者对。

\paragraph{连续时间状态转移}
设第 $k$ 次决策时点为 $\tau_k$，决策前状态为 $\nu^k$，选择
$a^k\in\mathcal A(\nu^k)$ 后得到 $v^k=v(\nu^k,a^k)$。
在下一次影响匹配池的到达或离开之前，池内人数保持 $v^k$。
令 $e_i$ 为第 $i$ 个单位向量，
\[
q(v)=\sum_{i=1}^d\lambda_i+\sum_{i=1}^d\mu_i v_i.
\]
由独立泊松到达和指数耐心的无记忆性，条件于当前已知历史及动作，
\begin{align}
\tau_{k+1}-\tau_k&\sim\operatorname{Exp}(q(v^k)),\tag{T1}\\
\Pr\{\nu^{k+1}=v^k+e_i\mid v^k\}&=\frac{\lambda_i}{q(v^k)},\quad i\in\mathcal T,\tag{T2}\\
\Pr\{\nu^{k+1}=v^k-e_i\mid v^k\}&=\frac{\mu_i v_i^k}{q(v^k)},\quad i\in\mathcal T:v_i^k>0.\tag{T3}
\end{align}
下一事件若为离开，该参与者先退出，不能在已离开后再选为匹配对象；
下一事件若为到达，新参与者进入后才作下一匹配决策。
并列事件的概率为零。池内离开事件由尚在池中的参与者产生。
时间按改变剩余池的事件进行索引。
取初始空池 $\nu^0=0,\tau_0=0,a^0=0$，与原文从首个到达开始构造系统一致。

\paragraph{信息约束与唯一优化目标}
令 $\mathcal H_k$ 为截至 $\tau_k$ 已观察到的到达、离开、状态与过去匹配动作
所生成的信息。允许独立内部随机数，但不允许未来信息。
一个可行政策 $\pi$ 在每个 $k$ 给出仅依赖 $\mathcal H_k$ 的动作分布，
其支持必须属于 $\mathcal A(\nu^k)$。记所有这样的政策为 $\Pi$。
优化模型为
\begin{equation}
\boxed{\quad\sup_{\pi\in\Pi}\ 
\liminf_{t\to\infty}\frac1t\,
\mathbb E^{\pi}\!\left[
\sum_{k:\,0<\tau_k\le t}g(a^k)
\right]\quad}\tag{DSM-R}
\end{equation}
并同时满足(A)、(S)、(R)、(T1)--(T3)及上述信息约束。
目标为原文第20--21页定义的政策收益上确界 $\sup_\pi r^\pi$。原文第10页说明，在其正则条件下可以选最优确定性平稳政策。
系统队长取非负整数，政策在无限时域上定义。

\end{document}
